\documentclass[10pt,a4paper]{amsart}
\usepackage[margin=19mm]{geometry}
\usepackage{amsmath,amssymb,mathtools}
\usepackage[scr=rsfs,cal=euler]{mathalfa}
\usepackage{xfrac}
\usepackage[unicode,hidelinks,bookmarks=false]{hyperref}
\newcommand{\ie}{i.e.\ }
\newcommand{\cf}{cf.\ }
\newcommand{\resp}{respectively\ }
\newcommand{\Subsection}{\subsection}
\providecommand{\nobreakdash}{\nobreak\hbox{-}}
\setlength{\emergencystretch}{2em}
\multlinegap0pt
\usepackage{tikz}
\usepackage{mathtools}
\usepackage{float}
\usepackage{xcolor}
\newtheorem{proposition}{Proposition}
\newtheorem{corollary}{Corollary}
\DeclareMathOperator{\Gr}{Gr}
\DeclareMathOperator{\Herm}{Herm}
\newcommand{\NJ}{\mathrm{NJ}}
\newcommand{\R}{\mathbb{R}}
\newcommand{\St}{\mathrm{St}}
\DeclareMathOperator{\U}{U}
\newcommand{\vol}{\mathrm{vol}}
\title{\bf Degree of tensor train varieties \\ via integral geometry}
\author{Andrea Rosana, Otto T.P. Schmidt}
\date{Source: arXiv:2606.11847v1 (CC BY 4.0)}
\begin{document}
\maketitle

\noindent\emph{Source and licence.} \bf Degree of tensor train varieties \\ via integral geometry. Version arXiv:2606.11847v1 (CC BY 4.0), distributed by arXiv under the Creative Commons Attribution 4.0 International licence, http://creativecommons.org/licenses/by/4.0/, as recorded in the arXiv licence metadata for this exact version and shown on the abstract page https://arxiv.org/abs/2606.11847v1. Full licence notice: \texttt{licenses/arxiv-2606.11847-CC-BY-4.0.txt}.\par\smallskip

\noindent\textit{Editorial scope.} Counted unit: the article's proposition prop:recursive\_coarea\_formula (source0.tex lines 1120--1128). The article's complete original proof of it is delivered with it (source0.tex lines 1129--1185, one \texttt{proof} environment of 57 lines). No mathematics is written, repaired or re-derived: the statement and the proof are the article's own lines, in the article's own notation, and no retained line is substituted or re-flowed.  Retained uncounted as the source-local context the proof needs: lines 321--323; lines 355--357; lines 818--825; lines 1043--1058; lines 1107--1115.  Nothing was omitted from any retained span, so the package carries no omission record.  No retained line contains a citation, so the wrapper carries no bibliography and no bibliography entry.  No retained line contains a figure environment, an image-inclusion command or drawing code, so no image is delivered and the package declares no asset dependency.  Editorial formatting only, in addition: an AMS \texttt{amsart} preamble that restates the article's own declarations and macros used by the retained lines, namely the environments \texttt{proposition}, \texttt{corollary} on separate counters, together with the standard packages those lines need.  The retained environments print in this document as Proposition 1, Corollary 1 in order of appearance, whereas the article prints the same items with the numbering of its own full document.  The complete native source of the article as distributed in the arXiv e-print for this exact version is bundled unchanged as \texttt{sources/202578/source0.tex}.
\begin{equation}\label{eq:vol_gr}
    \vol(\Gr(k,n)) = \frac{\vol(\St(k,n))}{\vol(\U(k))}\,,
\end{equation}

\begin{equation}\label{eq:coarea_general_clean}
    \int_X (h\circ\varphi)(x)\,\NJ(\varphi,x)\,d\vol_X(x) = \int_Y h(y) \vol_X(\varphi^{-1}(y))\,d\vol_Y(y)\, .
\end{equation}

\begin{proposition}\label{prop:fiber_volume_Psi_r}
The volume of the fiber of \(\Psi_r\) through $q_r=(M^r,[T_{r+1}]) \in Y_r$ is
\[
\vol\bigl(\Psi_r^{-1}(\Psi_r(q_r))\bigr)
=
\vol(\U(D_r))\,\sqrt{\det(I+B_r^*B_r)}. 
\]
\end{proposition}

\begin{corollary}\label{cor:NJ_Psi_r}
The normal Jacobian of the one-step map \(\Psi_r: Y_r \longrightarrow X_r\) is
\[
\NJ(\Psi_r,(M^r, [T_{r+1}]))
=
\det\bigl(\Sigma_{r+1}([T_{r+1}])\bigr)^{m_r}
\,J_r(\Psi_r(M^r, [T_{r+1}])),
\]
where we denote
\begin{equation}\label{eq:Sigma_r_definition}
\Sigma_r([T_r])
:=
T_r^{(r-1)}
(T_r^{(r-1)})^* \in \Herm_{>0}(D_{r-1}).
\end{equation}
\end{corollary}

\begin{equation}\label{eq:F_recursion}
    F_r(\Sigma)
:=
\det(\Sigma)^{m_r}
\int_{\St(D_r,n_r)}
F_{r-1}\!\left(
L_{M^r}(\Sigma)
\right)\,d\mu(M^r),
\end{equation}

\begin{proposition}\label{prop:recursive_coarea_formula}
The following recursive relation holds
\begin{align}\label{eq:integral_recursion}
\int_{X_r}F_{r-1}(\Sigma_r([T_r]))\,d\vol_{X_r}([T_r])
&=
\vol(\Gr(D_r, n_r))
\int_{X_{r+1}}F_r(\Sigma_{r+1}([T_{r+1}]))\,d\vol_{X_{r+1}}([T_{r+1}]).
\end{align}
\end{proposition}

\begin{proof}
We apply the smooth coarea formula to \(\Psi_r:Y_r\rightarrow X_r\) for the function $h_r:X_r\to \R$ defined as 
\begin{equation*}
h_r([T_r]):= \frac{F_{r-1}(\Sigma_r([T_r]))}{J_r([T_{r}])}\,.
\end{equation*}
By Proposition \ref{prop:fiber_volume_Psi_r}, the right-hand side of \eqref{eq:coarea_general_clean} reads
\begin{align*}
\int_{X_r}h_r([T_r])\,\vol(\Psi_r^{-1}([T_r]))\,d\vol_{X_r}([T_r]) &= \int_{X_r}
\frac{F_{r-1}(\Sigma_r([T_r]))}{ J_r([T_{r}])}
\cdot
\vol(\U(D_r))\,J_r([T_{r}])
\,d\vol_{X_r}([T_r]) \\
&= \vol(\U(D_r))
\int_{X_r}F_{r-1}(\Sigma_r([T_r]))\,d\vol_{X_r}([T_r]).
\end{align*}
Now consider the left-hand side of \eqref{eq:coarea_general_clean} given by
\begin{align*}
\int_{Y_r}
\NJ(\Psi_r,q_r)\,
h_r(\Psi_r(q_r))
\,dM^r\,d\vol_{X_{r+1}}([T_{r+1}]),
\end{align*}
where $q_r=(M^r,[T_{r+1}])\in Y_r$. Using Corollary \ref{cor:NJ_Psi_r}, and the recursion
formula
\[
\Sigma_r(\Psi_r(M^r,[T_{r+1}]))
=
\sum_{j_r=1}^{d_r}
M^r_{j_r}\Sigma_{r+1}([T_{r+1}])(M^r_{j_r})^* = L_{M^r}(\Sigma_{r+1}([T_{r+1}]),
\]
we obtain
\begin{align}
\notag &\int_{Y_r}
\det(\Sigma_{r+1}([T_{r+1}]))^{m_r}
J_r(\Psi_r(M^r,[T_{r+1}]))
\frac{
F_{r-1}\!\left(L_{M^r}(\Sigma_{r+1}([T_{r+1}]))
\right)
}{
J_r(\Psi_r(M^r,[T_{r+1}]))
}
\,dM^r\,d\vol_{X_{r+1}}([T_{r+1}]) \\ \label{eq:rhs_coarea}
&\qquad=
\int_{Y_r}
\det(\Sigma_{r+1}([T_{r+1}]))^{m_r}
F_{r-1}\!\left(L_{M^r}(\Sigma_{r+1}([T_{r+1}]))
\right)
\,dM^r\,d\vol_{X_{r+1}}([T_{r+1}]).
\end{align}
Recall that the Stiefel measure can be normalized as $dM^r=\vol(\St(D_r,n_r))\,d\mu(M^r)$. Using \eqref{eq:F_recursion}, \eqref{eq:rhs_coarea} can be written as
\begin{align*}
\vol(\St(D_r,D_{r-1}d_r))
\int_{X_{r+1}}
F_r(\Sigma_{r+1}([T_{r+1}]))\,d\vol_{X_{r+1}}([T_{r+1}]).
\end{align*}
Comparing the obtained left and right sides of \eqref{eq:coarea_general_clean}, from \eqref{eq:vol_gr} we get the recursive integral expression.
\end{proof}

\end{document}
